% !PRINTED 236
\weberpage{236}
%<<CONT
\[
\begin{aligned}
5\,\sqrt{c}^{\,4} - 4\,c_1\,\sqrt{c}^{\,3} + 3\,c_2\,\sqrt{c}^{\,2} - 2\,c_3\,\sqrt{c} + c_4 &= 0,\\
-10\,\sqrt{c}^{\,3} + 6\,c_1\,\sqrt{c}^{\,2} - 3\,c_2\,\sqrt{c} + c_3 &= 0,
\end{aligned}
\]
und diese Gleichungen zerfallen wegen des doppelten Zeichens von $\sqrt{c}$ in die vier
\[
\begin{aligned}
5\,c^2 + 3\,c_2\,c + c_4 &= 0,\\
4\,c_1\,c + 2\,c_3 &= 0,\\
10\,c + 3\,c_2 &= 0,\\
6\,c_1\,c + c_3 &= 0.
\end{aligned}
\]
Aus der zweiten und vierten dieser Gleichungen folgt
\[
c_1 = 0, \quad c_3 = 0,
\]
und dann aus der dritten und ersten
\[
c_2 = -\frac{10}{3}\,c, \quad c_4 = -5\,c^2,
\]
so dass also die Gleichung f\"ur $y$ die Gestalt erh\"alt
\weq{10}{y^5 - \frac{10}{3}\,c\,y^3 + 5\,c^2\,y + c_5 = 0.}
Diese Gleichung h\"angt noch von den beiden Parametern $c_5$ und $c$ ab; man kann sie auf eine Gleichung mit einem Parameter reduciren durch die Substitution
\weq{11}{y = \sqrt{\frac{c}{3}}\,z, \quad \gamma = c_5\,\sqrt{\frac{3^5}{c^5}},}
wodurch sie die einfache und elegante Form erh\"alt
\weq{12}{z^5 - 10\,z^3 + 45\,z + \gamma = 0.}
Dies ist die \gesperrt{Brioschi'sche} Normalform. Die Substitution (11) leidet aber an dem Uebelstande, dass sie noch eine Quadratwurzel enth\"alt; w\"ahrend in den Formeln (8) und (10) nur die zwei im vorigen Paragraphen besprochenen Quadratwurzeln vorkommen. Man kann aber auch nach einer Bemerkung von \gesperrt{Klein} auf rationalem Wege aus (10) zu einer Normalform kommen, die nur einen Parameter enth\"alt.

Setzt man z.\,B.
\weq{13}{y = \frac{3\,c_5}{c^2}\,\frac{1}{z}, \quad \gamma = \frac{9\,c_5^2}{c^5},}
so erh\"alt man aus (10) die Gleichung
\weq{14}{z^5 + 15\,z^4 - 10\,\gamma\,z^2 + 3\,\gamma^2 = 0.}
